\documentclass[10pt,a4paper]{article}
\usepackage[margin=2cm]{geometry}
\usepackage{multicol}
\usepackage{times}
\usepackage{graphicx}
\usepackage{float}
\usepackage{amsmath}
\usepackage{booktabs}
\usepackage{caption}

\title{\textbf{Segmentation of COVID-19 Infection Regions from Chest X-ray Images Using U-Net}}
\author{Dinh Thanh Hai -- 22BA13119}
\date{}

\begin{document}
\maketitle

\begin{multicols}{2}

\section*{Abstract}
Chest X-ray imaging is widely used to examine lung diseases such as COVID-19.
Accurate segmentation of lung and infection regions can help doctors analyze
disease severity more efficiently.
In this work, a deep learning model based on the U-Net architecture is applied
to segment lung regions from chest X-ray images.
Experiments are conducted on the COVID-QU-Ex dataset.
The qualitative results show that the model can accurately segment lung areas
and produces predictions that are close to the ground truth masks.

\section*{I. Introduction}
COVID-19 is a respiratory disease that mainly affects the lungs.
Chest X-ray images are commonly used because they are fast and low-cost.
However, manual annotation of lung or infection regions is time-consuming and
requires expert knowledge.

Recently, deep learning methods have achieved strong performance in medical image
segmentation.
Among them, U-Net is one of the most popular architectures.
Therefore, this study uses a U-Net model to automatically segment lung regions
from chest X-ray images.

\section*{II. Dataset Description}
The dataset used in this study is the COVID-QU-Ex dataset, which is publicly
available on Kaggle.
It contains chest X-ray images and corresponding segmentation masks.

\subsection*{A. Data Structure}
Each sample in the dataset consists of:
\begin{itemize}
    \item One grayscale chest X-ray image.
    \item One binary mask representing lung regions.
\end{itemize}

After preprocessing and matching, a total of 39,746 image–mask pairs were obtained.
The dataset was split into:
\begin{itemize}
    \item Training set: 31,796 images
    \item Validation set: 7,950 images
\end{itemize}

\subsection*{B. Sample Visualization}
Figure~1 shows several examples of X-ray images and their corresponding ground
truth masks.
The masks clearly highlight the lung regions.

\begin{figure}[H]
    \centering
    \includegraphics[width=\linewidth]{download-9.png}
    \caption{Examples of chest X-ray images and ground truth masks}
\end{figure}

\subsection*{C. Mask Overlay}
To better visualize the annotations, the masks are overlaid on the original
X-ray images.
As shown in Figure~2, the lung regions are highlighted in red.

\begin{figure}[H]
    \centering
    \includegraphics[width=\linewidth]{download-10.png}
    \caption{X-ray images with lung mask overlay}
\end{figure}

\section*{III. Methodology}

\subsection*{A. Model Architecture}
The U-Net architecture is used in this work.
It consists of an encoder and a decoder connected by skip connections.
The encoder extracts high-level features, while the decoder restores spatial
information.
Skip connections help preserve fine details in the segmentation output.

\subsection*{B. Implementation Details}
The model is implemented using Python and the PyTorch framework.
All images are resized to a fixed resolution and normalized to the range [0, 1].
Binary Cross-Entropy loss is used during training.
The Adam optimizer is applied to update the model parameters.

\section*{IV. Experimental Results}

\subsection*{A. Qualitative Results}
Figure~3 shows a comparison between the input X-ray image, the ground truth mask,
and the predicted mask generated by the U-Net model.
The predicted masks are visually close to the ground truth and accurately capture
the lung shapes.

\begin{figure}[H]
    \centering
    \includegraphics[width=\linewidth]{download-11.png}
    \caption{Comparison of X-ray image, ground truth mask, and predicted mask}
\end{figure}

\subsection*{B. Discussion}
The results demonstrate that the U-Net model can effectively segment lung regions
from chest X-ray images.
Most lung boundaries are well detected.
Some small errors appear near the edges, which may be caused by image noise or
variations in lung shape.

Compared to traditional image processing methods, the deep learning approach
provides more robust and accurate segmentation results.
The performance is comparable to other state-of-the-art segmentation methods
reported in recent studies.

\section*{V. Conclusion}
This work presents a deep learning approach for lung segmentation from chest
X-ray images using the U-Net architecture.
Experiments on the COVID-QU-Ex dataset show that the model can generate accurate
segmentation masks.
The results indicate that U-Net is a suitable model for medical image segmentation
tasks.

In future work, the performance can be further improved by using more advanced
architectures, additional data augmentation, or more training epochs.

\end{multicols}

\end{document}
